\section{Solutions for Chapter~\ref{ch:full_theories}}

\paragraph{Exercise~\ref{ex:full-relevant}.}
\phantomsection
\label{sol:ex:full-relevant}
Squaring the seven top dimensions gives
\[
 3^2+5^2+7^2+13^2+14^2+27^2+42^2
 =9+25+49+169+196+729+1764=2941.
\]
These entries follow the increasing dimensions in the chapter's
table. A scalar in a diagonal sector has equal grades on the two
sides, so any positive grade adds at least two to its dimension.
This explains why the relevant scalar list uses only top states.

For the spinning levels, each positive-weight top vector supplies
one grade-one descendant. It cannot vanish in a unitary module,
because $\|L_{-1}|h\rangle\|^2=2h\langle h|h\rangle>0$.
Thus $v_i(1)=a_i(1)-a_i(0)$. In the $d_4$ sector,
\[
 (h,\bar h)=\left(\frac5{13}+1,\frac5{13}\right),
 \qquad \Delta=\frac{23}{13},\quad s_{\rm spin}=1,
\]
and the multiplicity is $5(229-5)=1120$.
Interchanging the two grades gives another $1120$ states of
spin $-1$. Similarly, the $d_1$ top has weight $6/13$, giving
$\Delta=2(6/13)+1=25/13$ and $7(292-7)=1995$ primary
components for each spin sign.



\paragraph{Exercise~\ref{ex:full-E8states}.}
\phantomsection
\label{sol:ex:full-E8states}
Integer grading leaves exactly three pairs of weights at total
dimension two: $(1,1)$, $(2,0)$, and $(0,2)$.
At $(1,1)$ there are $248^2=61504$ scalar primaries.
The chiral character has $4124$ states at weight two.
Of these, $248$ are the independent derivatives $L_{-1}J^A$
of the currents, and one is $L_{-2}\one$, the stress state.
There is no $L_{-1}\one$ descendant. Subtracting leaves
$4124-248-1=3875$ chiral primaries. Each of the pairs
$(2,0)$ and $(0,2)$ contains one copy of this decomposition.
Thus the total is
\[
 61504+2(3875)+2(248)+2=69752.
\]
The last $498$ states are not Virasoro primaries.



\paragraph{Exercise~\ref{ex:full-tail}.}
\phantomsection
\label{sol:ex:full-tail}
For each omitted state, factor
$E^ke^{-bE}=e^{-b_0E}[E^ke^{-(b-b_0)E}]$
and bound the bracket by its supremum for $E\ge E_D$.
Positivity permits summation. With $a=b-b_0>0$, the
suprema are
\[
 \sup_{E\ge E_D}Ee^{-aE}=
 \begin{cases}E_De^{-aE_D},&aE_D\ge1,\\
 (ae)^{-1},&aE_D<1,\end{cases}
\]
and
\[
 \sup_{E\ge E_D}E^2e^{-aE}=
 \begin{cases}E_D^2e^{-aE_D},&aE_D\ge2,\\
 4/(a^2e^2),&aE_D<2.\end{cases}
\]
Heat capacity uses the variance
$\langle E^2\rangle-\langle E\rangle^2$, so one needs
control of the zeroth, first, and second weighted sums,
including their common denominator. A small unweighted
tail can still carry a substantial energy moment.



\paragraph{Exercise~\ref{ex:full-onevac}.}
\phantomsection
\label{sol:ex:full-onevac}
There are six dimension-zero states, which describes a
multiple-vacuum or disconnected situation instead of the
specified connected theory with a unique vacuum.
Division by six can destroy integral multiplicities and
does not construct identity fields, OPEs, or sewing maps.
A numerical rescaling of a partition function cannot repair
missing local-theory data.



\paragraph{Exercise~\ref{ex:full-symmetry}.}
\phantomsection
\label{sol:ex:full-symmetry}
The $36=6^2$ count is the diagonal theory's scalar
dimension-two subspace invariant under all its Verlinde lines.
The $1488=6\cdot248$ count is invariant under the selected
first-factor $\mathcal H_3$ action in the extended full theory.
The $61504=248^2$ count is the entire scalar dimension-two
space of that extended theory, with no symmetry restriction.
Different ambient Hilbert spaces and line actions are being used.

The invariant spaces themselves can also be checked.
In the diagonal theory, modular orthogonality makes the projector
in Chapter~\ref{ch:defects} select precisely
$U\otimes\overline U$. In the extension, use the line matrices
on $H'=U\oplus b\oplus(\C^2\otimes a)$.
The invariance conditions are
$\widehat\alpha\psi=\psi$ and $\widehat\rho\psi=d\psi$.
The $a$ multiplicity block has $\alpha$ eigenvalues
$\omega,\omega^2$, so it contributes no invariant vector.
The $b$ block survives $\alpha$ but has $\rho$ eigenvalue
$-d^{-1}\ne d$, so it also contributes none.
The entire $U$ block survives. Since the chosen first factor
does not act on the right movers, the full invariant space is
$U\otimes\overline{H'}$. Its scalar weight pair $(1,1)$
has $6\cdot248$ components. Imposing the independent right
action repeats the same argument and reduces this to $6^2$.



\paragraph{Exercise~\ref{ex:full-boost}.}
\phantomsection
\label{sol:ex:full-boost}
The two block products are
\[
 X^{\mathsf t}\eta=
 \begin{pmatrix}0&-\varepsilon\\
 \varepsilon^{\mathsf t}&0\end{pmatrix},\qquad
 \eta X=
 \begin{pmatrix}0&\varepsilon\\
 -\varepsilon^{\mathsf t}&0\end{pmatrix}.
\]
They sum to zero. Thus $e^{tX}$ preserves the quadratic
form $p^{\mathsf t}\eta p$, and its infinitesimal action is
the stated charge boost: differentiate $p(t)=e^{tX}p(0)$
at $t=0$ and read its two blocks. In particular,
\[
 \delta(p_L^2-p_R^2)
 =2p_L^{\mathsf t}\varepsilon p_R
  -2p_R^{\mathsf t}\varepsilon^{\mathsf t}p_L=0,
 \qquad
 \delta\Delta=2p_L^{\mathsf t}\varepsilon p_R.
\]
For $r=1$ and $\varepsilon=-1$,
$X^2=I$ and $e^{tX}=I\cosh t+X\sinh t$.
Starting at $(p_L,p_R)=(\sqrt2,0)$ yields
$(\sqrt2\cosh t,-\sqrt2\sinh t)$.
The spin remains $(p_L^2-p_R^2)/2=1$, although
$\bar h=\sinh^2t>0$ for $t\ne0$.
This is a statement about the charge contribution, with oscillator
and neutral labels held fixed. Establishing a local conformal
family uses the abelian current-current theorem in the chapter;
the matrix identity alone would not construct one or transport
the Haagerup defect junctions.


